\chapter{Referensi Tambahan dan Sumber Belajar}
\label{app:referensi}

\section{Referensi RPS}
\label{app:referensi-rps}

\begin{itemize}[leftmargin=*]
  \item \textbf{[1] Thain (2024)}: pengantar compiler modern dengan
        proyek B-Minor; sangat baik untuk scanning, parsing, dan AST.
  \item \textbf{[2] Mogensen (2010)}: teori ringkas dan padat;
        kuat untuk automata, parsing, dan type checking.
  \item \textbf{[3] Bergmann (2024)}: teori, alat, dan contoh;
        berguna untuk latihan terstruktur.
  \item \textbf{[4] Bornat (2008)}: panduan praktis menulis
        compiler; berguna untuk perspektif implementasi.
  \item \textbf{[5] Critchlow dan Eck (2024)}: dasar teori
        komputasi dan automata.
  \item \textbf{[6] Srikant}: catatan NPTEL; ringkas dan berorientasi
        ujian.
  \item \textbf{[7] MIT OCW 6.035}: kuliah dan tugas compiler.
  \item \textbf{[8] Stanford CS143}: slide dan handout parsing.
  \item \textbf{[9] Saarland}: struktur compiler.
  \item \textbf{[10] McGill CS520}: handout top-down parsing.
  \item \textbf{[11] Edinburgh}: compiling techniques.
  \item \textbf{[12] Flex Manual}: referensi lexer generator.
  \item \textbf{[13] GNU Bison Manual}: referensi parser generator.
  \item \textbf{[14] ANTLR 4}: dokumentasi parser generator.
  \item \textbf{[15] Wikibooks Compiler Construction}: ringkasan
        praktis.
  \item \textbf{[16] Wikibooks Formal Languages}: teori bahasa dan
        automata.
\end{itemize}

\section{Karya Lokal}
\label{app:referensi-lokal}

\begin{itemize}[leftmargin=*]
  \item Nystrom (2021), \emph{Crafting Interpreters}: implementasi
        interpreter dan VM, sangat baik untuk scanning dan parsing.
  \item Guntoro (2024), \emph{Teknik Kompilasi}: istilah Indonesia,
        penanganan kesalahan, dan tabel informasi.
  \item Juarna, \emph{Catatan Kuliah Teknik Kompilasi}: ringkasan
        leksikal, sintaksis, semantik, dan kode antara.
\end{itemize}

\section{Alat dan Sumber Daring}
\label{app:referensi-alat}

\begin{itemize}[leftmargin=*]
  \item Compiler Explorer (\code{godbolt.org}): melihat kode assembly
        dari berbagai compiler.
  \item regex101 (\code{regex101.com}): menguji ekspresi reguler.
  \item Flex manual (\code{westes.github.io/flex/manual/}).
  \item GNU Bison manual (\code{gnu.org/software/bison/manual/}).
  \item ANTLR Lab (\code{lab.antlr.org}).
  \item LLVM Kaleidoscope (\code{llvm.org/docs/tutorial/}).
  \item TinyCC (\code{bellard.org/tcc/}).
\end{itemize}

\section{Jalur Bacaan}
\label{app:referensi-jalur}

\begin{table}[H]
\centering
\caption{Jalur bacaan menurut minat}
\label{tab:app-jalur}
\begin{tabularx}{\textwidth}{@{}l X@{}}
\toprule
\textbf{Minat} & \textbf{Bacaan} \\
\midrule
Dasar & Thain Bab 1--5; Mogensen Bab 1--3 \\
Lexer & Thain Bab 3; Flex manual \\
Parser & Thain Bab 4--5; Bison manual; Stanford CS143 \\
Semantik & Thain Bab 7; Mogensen Bab 6 \\
Praktik & Nystrom; Juarna; Guntoro \\
Teori & Critchlow dan Eck; Wikibooks Formal Languages \\
\bottomrule
\end{tabularx}
\end{table}

\section{Peta Sitasi}
\label{app:referensi-peta}

\begin{table}[H]
\centering
\caption{Kunci bibliografi dan nomor RPS}
\label{tab:app-sitasi}
\begin{tabular}{@{}ll@{}}
\toprule
\textbf{Kunci} & \textbf{Nomor RPS} \\
\midrule
\code{thain2024} & [1] \\
\code{mogensen2010} & [2] \\
\code{bergmann2024} & [3] \\
\code{bornat2008} & [4] \\
\code{critchlow2024} & [5] \\
\code{srikant} & [6] \\
\code{mitocw2010} & [7] \\
\code{stanfordcs143} & [8] \\
\code{saarland2013} & [9] \\
\code{mcgill2026} & [10] \\
\code{edinburgh2024} & [11] \\
\code{flexmanual} & [12] \\
\code{bisonmanual} & [13] \\
\code{antlr4} & [14] \\
\code{wikibookscc} & [15] \\
\code{wikibookstfl} & [16] \\
\bottomrule
\end{tabular}
\end{table}
